\chapter{Perspectives d'amélioration}
\label{chap:perspectives}

Le MVP livré pose des fondations solides ; il ouvre néanmoins de nombreuses pistes
d'amélioration, dont plusieurs correspondent aux contraintes technologiques avancées du
sujet. Cette feuille de route est organisée par horizon.

\section{Court terme : finalisation des briques posées}

\begin{itemize}
  \item \textbf{Câblage de l'événementiel Kafka.} L'infrastructure (producteur configuré,
        topic \texttt{bizcore.events}, consommateur de notifications) est en place ; il
        reste à publier effectivement les événements de domaine (changement de statut d'une
        demande, génération de facture) pour rendre la chaîne pleinement \emph{event-driven}.
  \item \textbf{Activation du cache Redis.} Le gestionnaire de cache et les TTL sont
        définis ; l'ajout d'annotations \texttt{@Cacheable} / \texttt{@CacheEvict} sur les
        lectures fréquentes (règles métier, configuration tenant, catalogue) activerait le
        cache déjà préfixé par tenant.
  \item \textbf{Limitation de débit (rate limiting).} La dépendance Bucket4j est présente ;
        son intégration permettrait d'appliquer des quotas par tenant, prolongeant l'analogie
        de la bande passante et de la QoS.
  \item \textbf{Mesure de couverture de tests.} L'ajout de JaCoCo fournirait un indicateur
        chiffré de couverture.
    \item\textbf{Sur le plan de l'architecture: }il s'agirait d'étendre la couverture des cent cinquante-deux tests automatisés afin d'inclure des scénarios de charge et des tests d'intégration multitenant plus poussés, garantissant la robustesse de l'isolation par tenantId dans des conditions de concurrence élevée
    
    \item\textbf{La documentation OpenAPI: }pourrait également être enrichie d'exemples d'utilisation et de scénarios d'erreur, facilitant l'adoption de l'API par de nouveaux développeurs

    \item\textbf{le service d'analytique et la piste d'audit:} pourraient être complétés par des tableaux de bord de supervision en temps réel, exploitant les données déjà collectées pour offrir une visibilité immédiate sur l'activité de chaque tenant
    
\end{itemize}

\section{Moyen terme : robustesse et qualité de service}

\begin{itemize}
  \item \textbf{Isolation multi-tenant renforcée.} Passer de la discipline applicative à
        une barrière systématique au moyen d'un filtre Hibernate \texttt{@Filter} ou de la
        \emph{Row-Level Security} PostgreSQL, voire d'une isolation par schéma dédié pour
        les tenants premium.
  \item \textbf{QoS différenciée.} Introduire des niveaux de service inspirés de DiffServ
        (Bronze, Silver, Gold, Platinum) combinant priorité, latence cible et stratégie
        d'isolation.
  \item \textbf{Observabilité complète.} Exposer des métriques via Prometheus et un tableau
        de bord Grafana, et mettre en place des logs structurés agrégés.

  \item \textbf{Le modele generique}pourrait être étendu pour couvrir de nouveaux domaines métier au-delà du cas d'usage initial, validant ainsi la généricité réelle du Business Core sur des secteurs d'activité variés (santé, transport, logistique)

  \item \textbf{L'authentification JWT et le contrôle d'accès basé sur les rôles :} pourraient évoluer vers une gestion plus fine des permissions, avec la possibilité pour chaque tenant de définir ses propres rôles et politiques d'autorisation sans intervention sur le code du système
\end{itemize}

\section{Long terme : montée en charge et nouvelles contraintes}

\begin{itemize}
  \item \textbf{Architecture réactive.} Migration vers Spring WebFlux + R2DBC pour un
        modèle non bloquant, mieux adapté à une forte concurrence.
  \item \textbf{Géo-référencement natif.} Intégration de PostGIS et d'attributs de
        localisation pour les recherches géo-spatiales (\emph{geolocation aware}).
  \item \textbf{Stockage d'objets.} Adoption de MinIO pour les médias, au lieu de la simple
        URL actuellement stockée.
  \item \textbf{Sécurité renforcée.} Double authentification (2FA) et migration vers
        OAuth 2.0.
  \item \textbf{Front offline-first et international.} Transformation du frontend en PWA
        installable, mode hors-ligne, et internationalisation (i18n) web et mobile.
  \item \textbf{Application mobile native.} Client React Native partageant l'API.
  \item \textbf{Microservices.} Extraction progressive des modules en services déployés
        indépendamment, le découpage actuel ayant été pensé pour cette évolution.
  \item \textbf{Marketplace d'instances.} À terme, permettre à des entreprises de déployer
        leur propre configuration BCaaS sans développement, comme socle d'un écosystème
        d'instances métier.
  \item \textbf{ La grille de lecture protocolaire en cinq couches : }, actuellement formalisée et reliée au modèle OSI, pourrait servir de fondement à une véritable spécification ouverte du Business Core as a Service, documentée et partagée, permettant à d'autres équipes de contribuer à son évolution ou de l'adopter comme standard interne
\end{itemize}
